\documentclass[../main.tex]{subfiles}
\begin{document}

\begin{comment}
======================================
======================================
DÉVELOPPEMENT
======================================
======================================
\end{comment}

\chapter{Développement}

\section{Cadre théorique}
 
Texte explicatif appuyé par des références \cite{Exemple}.

\section{Méthodologie}
 
\subsection{Modèle utilisé}
 
Le calcul de la tension se fait selon :
\begin{equation}
    V = RI
\end{equation}
\begin{definitions}[où :]
    V & Tension aux bornes de la résistance [V] \\
    R & Résistance [$\Omega$]                   \\
    I & Courant traversant la résistance [A]    \\
\end{definitions}

\section{Résultats}

La figure~\ref{fig:exemple} illustre le comportement mesuré.
 
% Figure (commade définie dans app-utils.sty)
\quickfigure{Comportement du système mesuré}{exemple}{0.6}{exemple.jpg}
 
% Table (commade définie dans app-utils.sty)
\quicktable{NomTableau}{tab-name}{l c c}{
    \toprule
    Nom1 & Nom2 & Nom3 \\
    \midrule
    \texttt{NomText1} & Data1 & Data2 \\
    \texttt{NomText2} & Data3 & Data4 \\
    \texttt{NomText3} & Data5 & Data6 \\
    \midrule
    \texttt{NomText1} & Data1 & Data2 \\
    \texttt{NomText2} & Data3 & Data4 \\
    \midrule
    \texttt{NomText1} & Data1 & Data2 \\
    \texttt{NomText2} & Data3 & Data4 \\
    \bottomrule
}

% Les diagrammes UML sont lourd à rendre. Vous les faites soit hors de LaTeX,
% ou bien vous pré-compilez (comme mentionné dans le main.tex)
\section{Diagrammes UML}

% Use case diagram (commande définie dans app-diagrams.sty)
\begin{quickusecasediagram}[0.85]{Diagramme de cas d'utilisation}{usecase-exemple2}
 
    % Cadre du système
    \umlsystem{-1.8}{-5.5}{8.5}{4.5}{QtDoom}
 
    % Acteur
    \umlactor{Utilisateur}{-4.5}{0}
 
    % Cas d'utilisation
    \umlusecase{uc1}{3}{3.0}{Entrer son nom}
    \umlusecase{uc2}{3}{1.5}{Entrer mot de passe}
    \umlusecase{uc3}{3}{0.0}{Valider connexion}
    \umlusecase{uc4}{3}{-1.5}{Afficher profil}
    \umlusecase{uc5}{3}{-3.0}{Se déconnecter}
    \umlusecase{uc6}{3}{-4.5}{Modifier profil}
 
    % Associations d'acteurs et de cas d'utilisation
    \draw[umlassoc] (-4.1, 0.9) -- (uc1.west);
    \draw[umlassoc] (-4.1, 0.9) -- (uc2.west);
    \draw[umlassoc] (-4.1, 0.9) -- (uc3.west);
    \draw[umlassoc] (-4.1, 0.9) -- (uc4.west);
    \draw[umlassoc] (-4.1, 0.9) -- (uc5.west);
    \draw[umlassoc] (-4.1, 0.9) -- (uc6.west);
 
    % Relation include
    \umlrelation[right]{include}{uc6}{uc5}{}
    \draw[umlextend] (uc3.east) 
                     -- ++(1, 0) 
                     -- node[right, font=\scriptsize]{<<extend>>} ++(0, -1.5) 
                     -- (uc4.east);
 
\end{quickusecasediagram}

% Class diagram (commande définie dans app-diagrams.sty)
\begin{quickclassdiagram}[0.7]{Diagramme de classes}{classes-exemple3}
 
    % Classe de librairies externes 
    \node[umlclass] (QWidget) at (0, 0) {QWidget};
 
    % Classe crée sois même
    \node[umlclassnode] (APP) at (0, -6) {
        \begin{umlclassbox}[6cm]{APP}
            \umlattr{- acceuil : Acceuil*}
            \umlattr{- compte : Compte*}
            \umlattrbreak
            \umlmethod{+ APP(parent : QWidget*)}
            \umlmethod{+ deconnexion()}
            \umlmethod{+ loadCompte(username : QString)}
        \end{umlclassbox}
    };
 
    % Classe Acceuil
    \node[umlclassnode] (ACC) at (-7, -6) {
        \begin{umlclassbox}[5.5cm]{Acceuil}
            \umlattr{- profilPage : Profil*}
            \umlattr{- statsPage : Stats*}
            \umlattrbreak
            \umlmethod{+ Acceuil(parent : QWidget*)}
            \umlmethod{+ deconnexion()}
            \umlmethod{+ openProfilPage()}
        \end{umlclassbox}
    };
 
    % Classe Compte
    \node[umlclassnode] (COMPTE) at (7, -6) {
        \begin{umlclassbox}[5.5cm]{Compte}
            \umlattr{- username : QString}
            \umlattr{- nip : QString}
            \umlattrbreak
            \umlmethod{+ Compte(parent : QWidget*)}
            \umlmethod{+ changeUser(...)}
        \end{umlclassbox}
    };
 
    % Héritages
    \draw[umlheritage] (APP.north) -| (QWidget.south);
    \draw[umlheritage]
        (ACC.north) -- ++(0, 0.5) 
                    -- ++(3.6, 0) 
                    -- ++(0, 3.25) 
                    -- (QWidget.west);
    \draw[umlheritage]
        (COMPTE.north) -- ++(0, 0.5) 
                       -- ++(-3.6, 0) 
                       -- ++(0, 3.57) 
                       -- (QWidget.east);
 
    % Compositions
    \umlrelation{comp}{APP.west}{ACC.east}{}
    \umlrelation{comp}{APP.east}{COMPTE.west}{}
 
\end{quickclassdiagram}

\end{document}